% 第 3 章 §3.2（印刷页 118–127）
% 对应 工程文件/work/en/ch03.txt 第 475–1000 行：§3.2 标题至定理 3.2.8 后的讨论。
% 注意：第 951–1000 行虽在印刷页 127，仍属 §3.2；§3.3 自第 1001 行起（见 ch03c.tex 文件头）。
% 编号公式仅 (3.2)(3.3)(3.4)；无插图，两个交换图用 tikzcd 重绘；脚注 8 条（原书 6–13）。

\section{二维规范理论}

自从 Witten [31] 用路径积分中的局部化论证显式计算了这一理论、令数学界大吃一惊以来，这个例子就一直是人们的心头好。其结果给出曲面 $\Sigma$ 上平坦 $G$-联络的模空间 $\Phi(\Sigma;G)$ 上一切积分的 TQFT 计算，所用的只是一个显式算出的Frobenius 代数。为方便起见，设 $G$ 为单且单连通；此时 $\Phi(\Sigma;G)$ 的维数为 $2d=2\dim G\cdot(g-1)$。我们暂且无视一个令人不便的事实：$\Phi$ 通常含有奇点，对应于和乐表示可约的丛。

$\Sigma\times\Phi(\Sigma;G)$ 上的万有平坦丛由一个到分类空间 $BG$ 的映射 $u$ 分类（$u$ 在相差同伦的意义下确定）。回忆 $H^*(BG;\C)$ 是多项式环，同构于 $(\operatorname{Sym}\mathfrak{g}^*)^G\cong(\operatorname{Sym}\mathfrak{t}^*)^W$，生成元 $\varphi_2,\dots,\varphi_\ell$ 的上同调次数\footnote{对称代数的次数（作为 $\mathfrak{g}$ 上的多项式）是上同调次数的一半。}为 $2m_i+2$，其中 $m_i$ 是该群的指数。（例如，$\mathrm{SU}(n)$ 的指数为 $1,2,\dots,n-1$。）取 $\Sigma$ 的同调的一组基中的元素 $\beta$，斜积 $\beta\backslash u^*\varphi_i$（即 $u^*\varphi_i$ 在 $\beta$ 上的积分）给出次数为 $2m_i+2$、$2m_i+1$ 与 $2m_i$ 的类。Atiyah--Bott [2] 证明：只要 $\Phi(\Sigma;G)$ 光滑，$H^*\Phi(\Sigma;G)$ 就由这些类生成。\footnote{光滑性只可能对变体 $\Phi_c(\Sigma;G)$ 出现；在奇异情形，正确的陈述涉及等变上同调，而不是商空间的上同调。}

次数为 $2m_i$ 的类 $[\Sigma]\backslash u^*\varphi_i$ 在积分中最「有趣」，因为它们不容易化归为初等计算。例如，由对应于 $\mathfrak{g}$ 上不变二次型 $\langle\,,\rangle$ 的二次 Casimir 元 $\varphi_2$，我们得到 §3.1 中提到的辛形式 $\omega$。（二次型有一个优选的归一化，称为基本的归一化：在其中，$\mathfrak{g}$ 中最短的非零 $\log(1)$ 的长度平方为 $2$。）对主丛 $P$ 处的两个切向量 $\delta A,\delta B$（用 $\Sigma$ 上取值于 $\mathrm{ad}_P$ 的一形式表示），缩并 $\langle\delta A,\delta B\rangle$ 自然是 $\Sigma$ 上的 $2$-形式，于是
\[
\omega(\delta A,\delta B)=\int_\Sigma\langle\delta A,\delta B\rangle.
\]

\begin{annote}
底本此句印作「the shortest non-zero $\log(1)$ in $\mathfrak{g}$ has square length 2」，字面不可解。按上下文（basic 归一化），它应指最短的非零「余根」（coroot）：在该归一化下最短非零余根的长度平方为 $2$。译文照排底本字面。[译者按（上下文推断，无外部依据）：$\log(1)$ 当为 coroot 的排印错误]
\end{annote}

在 $\mathrm{SU}(2)$ 的特殊情形，Witten 的公式如下。取一个多项式 $Q$。对每个 $k\in\Z_+$，令 $\xi(k,t)$ 为（关于 $t$ 的）形式幂级数，它表示函数关于 $\xi$ 的唯一临界点（该函数还依赖于附加参数 $h,k,t$），并令
\begin{equation}
F(\xi;h,k,t):=\tfrac12(\xi-k)^2+t\cdot\frac{h}{2\pi^2}\cdot Q(\pi\xi/h),
\tag{3.2}\label{eq:3.2}
\end{equation}
其中 $t$ 被当作形式变量，故 $\xi(k,t)$ 由在极小点 $\xi(k,0)=k$ 附近的 $t$-级数展开算出。Witten 公式的一个特殊情形则说
\[
\begin{aligned}
&\int_{\Phi(\Sigma;G)}\exp\bigl\{h\omega+t\cdot[\Sigma]\backslash Q(u^*\varphi_2)\bigr\}\\
&\qquad=\#Z(G)\cdot h^{3g-3}\mathrm{vol}(G)^{2g-2}
\sum_{k>0}\left[\frac{1+\frac{t}{2h}Q''(\xi(k,t))}{\xi(t,k)^2}\right]^{g-1}.
\end{aligned}
\]

当然，对 $\mathrm{SU}(2)$ 有 $\#Z(G)=2$。伸缩因子 $h^{3g-3}$ 其实应被体积吸收，因为我们应当把基本度量按 $h$ 重新缩放。（对 $\mathrm{SU}(2)$，在基本度量下 $\mathrm{vol}(G)=1/\pi\sqrt2$；重新缩放后为 $h^{3/2}/\pi\sqrt2$。）我们在右端认出了半单（但无穷维）Frobenius 代数 $\bigoplus_k\C P_k$ 的亏格 $g$ 配分函数，其中对每个正整数有一个投影算子 $P_k$，其迹为
\begin{equation}
\mathrm{vol}_h(G)^{-2}\,\frac{\xi(t,k)^2}{1+\tfrac{t}{2h}Q''(\xi(k,t))}.
\tag{3.3}\label{eq:3.3}
\end{equation}
当 $h=1$ 且 $t=0$ 时，它们化归为 §3.1.7 的值，理当如此。

要得到最一般的积分，取若干多项式 $Q_i$，把它们与独立的形式变量 $t_i$ 耦合；对类似公式取 $\partial^n/\partial_1\cdots\partial_n$ 导数，便算出
\[
\int_{\Phi(\Sigma;G)}\exp[h\omega]\wedge\bigl([\Sigma]\backslash Q_1(u^*\varphi_2)\bigr)\cdots\wedge\bigl([\Sigma]\backslash Q_n(u^*\varphi_2)\bigr).
\]

\begin{remark}
对单群 $G$，最一般的公式还要纳入 $H^*\Phi(\Sigma;G)$ 的其余生成元。$Q$ 将是李代数 $\mathfrak{g}$ 上的不变多项式，它由在嘉当子代数 $\mathfrak{t}$ 上的限制决定。整数 $k$ 的角色由 $G$ 的最高权（用外尔向量 $\rho$ 移位）接替。Frobenius结构常数 (3.3) 中的分子是 (3.2) 的多变量函数 $F$ 的 Hessian 行列式；$\xi(t,k)^2$ 则换成 $\xi(t,k)$ 在 $\mathfrak{g}$ 中伴随轨道的体积。
\end{remark}

\subsection{K-理论中的解释}

解释并理解 Witten 的公式并非易事。例如，上面定义的示性类其实并不真正落在模空间 $\Phi$ 上——后者是奇异的。\footnote{与原 $\Phi$ 密切相关的、值得注意的一族光滑 $\Phi_c$ 来自 $G=\mathrm{SU}(n)$，其中 $c$ 是 $n$ 阶中心元。}事实证明，这些空间确有优选的（轨形）消解；这是解释这些公式的一条途径。这里我要讨论一种用（扭转）K-理论的解释，它使该场论可以有拓扑的描述与计算。\footnote{基于 Freed--Hopkins--Teleman [13] 与 Teleman--Woodward [28] 两篇论文。}

先回忆复（射影）流形 $X$ 上全纯向量丛 $E\to X$ 的 Hirzebruch--Riemann--Roch 定理：
\begin{equation}
\operatorname{Ind}(X;E):=\sum_k(-1)^k\dim H^k(X;\mathcal{O}(E))=\int_X\operatorname{ch}(E)\operatorname{Td}(X),
\tag{3.4}\label{eq:3.4}
\end{equation}
其中 $H^k$ 表示以全纯截面层 $\mathcal{O}(E)$ 为系数的 $k$ 次上同调（等于以 $E$ 为系数的 Dolbeault 上同调），而 Chern 特征 $\operatorname{ch}(E)$ 与 Todd 类 $\operatorname{Td}(X)$ 分别是 $E$ 与 $X$ 的切丛的示性类。具体地说，
\[
\operatorname{ch}(E)=\operatorname{Tr}\exp\Bigl(\frac{-R_E}{2\pi i}\Bigr),\qquad
\operatorname{Td}(X)=\det\frac{-R_{TX}/2\pi i}{1-\exp(-R_{TX}/2\pi i)},
\]
其中 $R$ 表示相应的丛 $E$ 或切丛 $TX$ 上某个 Hermite 联络的曲率形式。若 $E$ 是线丛 $L_i$ 的直和，另有定义 $\operatorname{ch}(E)=\sum\exp c_1(L_i)$；对 $\operatorname{Td}$ 有类似（但是乘性的）手法。由拓扑 K-理论的分解原理，只要顾及直和下的行为
\[
\operatorname{ch}(E_1\oplus E_2)=\operatorname{ch}(E_1)+\operatorname{ch}(E_2),\qquad
\operatorname{Td}(E_1\oplus E_2)=\operatorname{Td}(E_1)\cdot\operatorname{Td}(E_2),
\]
线丛这一特殊情形就足以定义 $\operatorname{ch}$ 与 $\operatorname{Td}$。

其实，从 $X$ 上全纯向量丛到 $\Z$ 的加性映射 $\operatorname{Ind}$ 可以延拓到一切拓扑丛，从而由复拓扑向量丛的 Grothendieck 群 $K^0(X)$ 给出一个线性指标映射 $\operatorname{Ind}\colon K(X)\to\Z$。（后一个群称为 $X$ 的偶拓扑 K-理论，是某个非寻常的、即广义的上同调理论的零次部分。）公式 (3.4) 把指标分解为：先经环同态 $\operatorname{ch}\colon K^0(X)\to H^{\mathrm{ev}}(X)$，再对 Todd 类作积分。

\begin{remark}
$\operatorname{ch}$ 给出代数同构 $K^0(X)\otimes\Q\cong H^{\mathrm{ev}}(X;\Q)$；但把指标映射当作迹，会在 $K^0(X)\otimes\Q$ 上定义出一个与 $H^{\mathrm{ev}}(X;\Q)$ 上不同的Frobenius 代数结构。
\end{remark}

\subsection{从指标到积分}

暂设 $E$ 是 Chern 类为 $\omega$ 的线丛。有一种办法把 $X$ 的辛体积作为指标的渐近提取出来：
\[
\operatorname{Ind}(X;E^{\otimes n})=n^{\dim X}\int_X\exp(\omega)+O(n^{\dim X-1}),
\]
这使得 Todd 类在首项中消失。对任意向量丛 $E$ 有类似的技巧，但不用张量幂，而要用第 $n$ 个 Adams 运算 $\psi^n$：对线丛 $L$ 定义 $\psi^nL=L^{\otimes n}$，并要求加性 $\psi^n(E\oplus F)=\psi^n(E)+\psi^n(F)$。（由此，分解原理唯一地确定了 $\psi^n\colon K^0(X)\to K^0(X)$。）在 Chern 特征下，$\psi^n$ 对应 $H^{2k}$ 上乘以 $n^k$。$\psi^n$ 还有一个用外幂表示的表达式，但其中带符号，故只在 Grothendieck 群 $K^0(X)$ 中有效，而在向量丛范畴中没有意义。无论如何，若把 $\psi^n$ 的计算视为已知，则由 $\psi^n$ 的线性可见，$X$ 上的积分可作为指标的渐近找回：
\[
\operatorname{Ind}(X;\psi^nE)=n^{\dim X}\int_X\operatorname{ch}(E)+O(n^{\dim X-1}).
\]

\subsection{$\Phi(\Sigma;G)$ 上的指标公式}

这段冗长铺垫的要点如下：$\Phi(\Sigma;G)$ 上的积分公式可以从指标公式找回。然而指标公式与积分不同，它来自一个真正的 TQFT，其基础是一个有限维Frobenius 代数。此外，指标公式可以纯粹几何地、由平坦联络模空间的对应图导出，正如有限群规范理论的情形。

一句哲学式的评论或许有帮助，或许没有——指标公式代表的是 $G$ 的圈群的一种规范理论版本。圈群的表示论令人想起紧李群的表示论，但其中带有一个基本的离散参数，即级（level）；而在固定的级上只有有限多个表示。于是固定级上的理论带有有限群表示论的某些特征。

最容易的故事与辛体积的 K-理论类比有关。当 $h$ 为整数时，$\exp(h\omega)$ 是 $\Phi(\Sigma,G)$ 上某个线丛 $\mathcal{O}(h)$ 的 Chern 特征。$\Sigma$ 上复结构的选择给出簇 $\Phi(\Sigma,G)$ 上的复解析（事实上是射影代数）结构：这是 Narasimhan 与 Seshadri 的一个深刻定理，它把后者等同于 $\Sigma$ 上（多）稳定全纯 $G$-丛的模空间。线丛 $\mathcal{O}(h)$ 其实容许全纯结构；当群 $G$ 单时，这种结构唯一。\footnote{对 $\U(1)$，$\Phi(\Sigma,G)$ 是 $\Sigma$ 的雅可比簇，全纯线丛可连续族地变化。}

$\mathcal{O}(h)$ 的指标有一个非常好的解释，这要归功于下述结果。

\begin{theorem}[Kumar--Narasimhan]
若 $h\ge0$，则上同调 $H^{>0}(F(\Sigma,G);\mathcal{O}(h))$ 为零。
\end{theorem}

\begin{annote}
底本在此处及后文若干处以 $F(\Sigma,G)$ 记同一个平坦联络模空间，与前文的 $\Phi(\Sigma;G)$ 混用；译文照排底本记号。[来源：底本 p.121 定理 3.2.3 与 p.122 Verlinde 公式]
\end{annote}

这些指标因而给出向量空间的维数。这些空间即共形块，已被大量研究。其维数的公式由 E. Verlinde 对 $\mathrm{SU}(2)$ 猜想：
\[
\dim H^0(F(\Sigma,G);\mathcal{O}(h))=(2h+4)^{g-1}\sum_{k=1}^{h+1}\Bigl(2\sin\frac{k\pi}{h+2}\Bigr)^{2-2g},
\]
而对一般的 $G$，它已在众多作者的工作中得到证明。Witten 的辛体积公式可由 Verlinde 公式的渐近得到。这里不复述推导，一般情形见 [28, §5]。我的目标毋宁是解释：为什么这些数（以及它们对向量丛指标的推广）受一个二维 TQFT 支配。

\subsection{Verlinde 环}

支配这些线丛指标的二维 TQFT，是一个半单的Frobenius $\Z$-代数，称为 Verlinde 环 $V(G;h)$。（到向量丛的推广则不那么为人所知；[27] 中已有所指示。）我这里只想指出：支配该理论的Frobenius 代数中的投影算子，与 $G$ 的光滑圈群 $LG$ 的投射、正能量表示自然对应，其投射余圈由线丛的 Chern 类 $h$ 决定。

圈群表示与 Verlinde 环之间确实有更深的联系：例如 Verlinde 环就是上述表示范畴的 Grothendieck K-群，该范畴上有一个张量结构，即融合积，它诱导 Verlinde 环中的乘法；有一个特殊的真空表示充当单位，而一般表示的Frobenius迹提取真空表示的重数。你会注意到这与有限群规范理论在形式上的相似：那里的「真空表示」就是平凡表示。不过我们跑题了。

\begin{remark}
从 TQFT 的观点看，我们赋给曲面的数是向量空间 $H^0$ 的维数。记住空间本身，便提示了三维 TQFT 的可能性：赋给曲面的是向量空间，赋给三维流形的是数。确实，这个 Chern--Simons 理论已被严格构造 [25]。这一主题通向共形场论、圈群表示与三维拓扑之间的深刻联系。
\end{remark}

\subsection{扭转 K-理论速成}

要描述支配 Verlinde 数及其（下面要讲的）推广的 TQFT，我们还需要一味原料。严格地构造扭转 K-理论，尤其是等变版本，本身就得花一门课；但想法相当简单。著名的 Serre--Swan 定理断言：紧 Hausdorff 空间 $X$ 上的向量丛恰是连续函数环 $C^0(X)$ 上的投射模。这给出 $K^0(X)$ 的一个纯代数定义：即相应的投射模的 Grothendieck 群。

当紧群 $G$ 作用在 $X$ 上时，可定义等变 K-群 $K^0_G(X)$，即带有 $G$-作用提升的向量丛的 Grothendieck 群。

这里同样有代数描述。可以作交叉积代数 $G\ltimes C^0(X)$：它是 $G\times X$ 上的函数，在 $X$ 上逐点相乘、在 $G$ 上作卷积，并用到 $G$ 在 $X$ 上的交织作用。例如 $G$ 有限时，$G\ltimes C^0(X)$ 中的元素可写成和 $\sum_{g\in G}g\cdot\varphi_g$，其中 $\varphi_g\in C^0(X)$，而乘法由
\[
\Bigl(\sum_{g\in G}g\cdot\varphi_g\Bigr)\cdot\Bigl(\sum_{h\in G}h\cdot\psi_h\Bigr)
=\sum_{k\in G}k\cdot\sum_{gh=k}(h^{-1})^*\varphi_g\cdot\psi_h
\]
给出，其中对 $u\in G$，$(u^*\varphi)(x):=\varphi(u^{-1}(x))$。换言之，把一个群元素移过去时，要对函数作用一下。

在这一语境中，Serre--Swan 定理的等变版本把 $K^0_G(X)$ 等同于 $G\ltimes C^0(X)$ 上投射模的 Grothendieck 群。

现在设想 $X$ 上的一个代数丛，局部同构于常值丛 $\C$。其实这样的丛只能是乘积丛；但这是因为对代数丛而言「同构」是错误的概念，应当代之以 Morita 等价。最简单的例子是矩阵代数丛。这样的丛不必整体平凡，甚至不必 Morita 等价于平凡丛。确实，给定投射向量丛 $\PP\to X$，相伴的矩阵代数丛 $\operatorname{End}(\PP)$ 是良定义的：这是因为对向量空间 $E$，$\operatorname{End}(E)$ 仅由射影空间 $\PP E$ 就典范地确定。然而 $\operatorname{End}(\PP)$ 不必是某个整体定义的向量丛 $E\to X$ 的自同态代数。确实，其障碍恰是把 $\PP$ 提升为向量丛 $E$（使 $\PP\cong\PP E$）的障碍。由短正合列
\[
1\to\GL(1)\to\GL(n)\to\operatorname{PGL}(n)\to1
\]
得到长正合列的片段
\[
\cdots\to H^1(X;\GL(1))\to H^1(X;\GL(n))\to H^1(X;\operatorname{PGL}(n))\to H^2(X;\GL(1)),
\]
它把障碍定位在 $H^2(X;\GL(1))$ 中。我们身处拓扑之中、用的是连续系数，而指数映射表明这与 $H^3(X;\Z)$ 相同。文献中称之为我们的投射丛所定义的 gerbe 的 Dixmier--Douady 类。

所得类都是 $n$-挠的；你们中间熟练的人会通过对比正合列
\[
1\to\mu_n\to\mathrm{SL}(n)\to\operatorname{PGL}(n)\to1
\]
注意到这一点，该列把障碍放在 $H^2(X;\mu_n)$ 中。不过，这一构造有一个好的无穷维版本，用的是投射 Hilbert 丛；在那里与 $\mathrm{SL}(n)$ 作比较的技巧失效了；而事实上可以证明，$H^3(X;\Z)$ 中的任何类都可由一个 Hilbert gerbe 实现，且在某种等价意义下唯一。一个技术性调整是：在 Hilbert 丛的故事里，矩阵代数丛的良好类比是紧自同态丛。

总而言之，对每个类 $[\tau]\in H^3(X;\Z)$，可定义扭转 K-群 ${}^\tau K^0(X)$，即由 $\tau$ 所定义的矩阵代数丛的截面上的投射模的 Grothendieck 群。除开一些已由多种方式解决的技术困难 [3, 14]，这个故事逐字延拓到带紧群作用的 $X$：$H^3_G(X;\Z)$ 中的一个类定义 $X$ 上一个 $G$-等变矩阵代数丛。（除非该类恰是挠类，这些丛是无穷维的。）从精神上说，我们定义扭转 K-群 ${}^\tau K^0_G(X)$ 为有限生成投射模的 Grothendieck 群。\footnote{再者，无穷维情形下 K-理论的定义还需一处技术性改动，因为紧算子构成非单位代数：添入一个单位，并限制到秩为 $0$ 的模；见 [3]。}

\begin{remark}[$\tau K$ 的 Chern 特征计算]
前面提到，一个空间的有理 K-理论相当容易由有理上同调算出。扭转 K-理论在有理系数下也相当容易计算，如下（遗憾的是，等变版本要麻烦得多）：
\begin{enumerate}[label=(\roman*),leftmargin=2.6em,itemsep=2pt]
\item 有一个奇 K-群 $K^1(X)$；一个偷懒的定义是 $K^1(X)=K^0(S^1\times X,X)$。这是 $K^{-1}$ 的定义，而我们在这里内置了 Bott 周期性定理 $K^i\cong K^{i+2}$。
\item 如前所述，对合理的紧空间 $X$，Chern 特征给出同构 $\operatorname{ch}\colon K^0(X)\otimes\Q\to H^{\mathrm{ev}}(X;\Q)$。相应地也有同构 $\operatorname{ch}\colon K^1(X)\otimes\Q\to H^{\mathrm{odd}}(X;\Q)$。今取 $(C^*(X;\Q),\delta)$ 为 $X$ 上有理上链构成的代数，微分记作 $\delta$。（若肯改用实系数，可以用 de Rham 的微分形式。）再取表示扭转类的 $3$-余圈 $\tau$。在这一语言中描述扭转 K-理论很容易：空间 ${}^\tau K^{0|1}$ 经由一个扭转 Chern 特征，同构于 $C^*(X;\Q)$ 在修改后的微分 $\delta+\tau\wedge$ 下的奇、偶上同调。（注意 $\delta\tau=\tau\wedge\tau=0$，这验证了 $(\delta+\tau\wedge)^2=0$；所以这是一个复形。）
\end{enumerate}
可惜，这个容易的模型对等变扭转 K-理论没有帮助；见 [12] 中关于非局部化 Chern 特征的更复杂的故事。
\end{remark}

\subsection{扭转的 $K_G(G)$}

设 $G$ 是紧、单、单连通的李群。则 K-理论的扭转群 $H^3_G(G;\Z)$ 典范同构于 $\Z$。（有一个优选的生成元，给出李代数上正定的基本二次型。）原来，级（即圈群表示的投射余圈）也由整数参数化。\footnote{投射 $LG$-余圈在拓扑上总由 $H^3_G(G;\Z)$ 分类。}[13] 的关键定理是Frobenius环的同构 ${}^\tau K^{\dim G}_G(G)\cong V(G;h)$。扭转 $\tau$ 为 $h+c$，其中级的移位 $c$（对偶 Coxeter 数）依群而定；对 $\mathrm{SU}(n)$ 它等于 $n$。在 $\mathrm{SU}(2)$ 的 Verlinde 公式中，$c=2$ 已经露过面了。

我们已暗示过 $V(G;h)$ 上的乘法；${}^\tau K^{\dim G}_G(G)$ 上的乘法是 Pontrjagin 积，由乘法 $G\times G\to G$ 诱导。若想到特征标的卷积，便又会想起有限群规范理论。这一次，我们面对的不是 $G$ 上的 $G$-不变函数，而是 $G$-等变向量丛（更确切地说，被 $\tau$ 扭转的投射丛）。用向量丛——取值于向量空间的函数——取代复值函数，这是范畴化的一个例子，也暗示我们所描述的 Verlinde 理论其实是一个伪装的三维理论。

${}^\tau K^{\dim G}_G(G)$ 上的Frobenius迹则不那么显然；见下文定理 3.2.7(iii)。在陈述它之前，回忆：圆周上平坦 $G$-联络的群胚模去规范等价，等价于商群胚 $G/G$，其中 $G$ 通过共轭作用在自身上。紧定向曲面 $\Sigma$ 带入边界 $\partial_-\Sigma$ 与出边界 $\partial_+\Sigma$，定义了平坦联络群胚之间的一个对应
\begin{center}
\begin{tikzcd}[row sep=1.8em, column sep=3.2em]
& \tilde\Phi(\Sigma;G) \arrow[dl, "\pi_+"'] \arrow[dr, "\pi_-"] & \\
\tilde\Phi(\partial^+\Sigma;G)\cong(G/G)^q & & (G/G)^p\cong\tilde\Phi(\partial^-\Sigma;G).
\end{tikzcd}
\end{center}

叠的语言把群胚结构与拓扑、微分性质结合起来。粗略地说，上面的叠 $\tilde\Phi$ 与底层的商模空间之别，在于前者记住了丛的自同构。不喜欢叠的人则必须等变地工作：在 $\Sigma$ 的（每个连通分支中）选一个基点、在每个边界分支中选一个基点，并选取从 $\Sigma$ 的基点到其余各基点的道路。于是每个叠都呈现为许多个 $G$ 的乘积——每条道路给一个平行移动，曲面或圆周的 $\pi_1$ 的每个生成圈给一个和乐——再除以 $G$ 在所有基点上的同时作用。这样做的缺点是：你接着必须证明下文定理 3.2.7 中的运算与基点、道路的选取无关。若把 K-理论与这些运算定义在叠及可表态射上，就可避开这一点。

\begin{exercise}
\begin{enumerate}[label=(\roman*),leftmargin=2.6em,itemsep=2pt]
\item 从某个呈现出发并尽可能化简，证明由带两个入圆周与一个出圆周的裤子所诱导的对应图是
\begin{center}
\begin{tikzcd}[row sep=1.8em, column sep=3.4em]
& (G\times G)/G^{\mathrm{diag}} \arrow[dl, "m"'] \arrow[dr, "\Id/\mathrm{diag}"] & \\
G/G & & G/G\times G/G,
\end{tikzcd}
\end{center}
其中所有群都通过共轭作用，$G^{\mathrm{diag}}$ 表示对角子群，$m$ 是群上的乘法映射。于是该运算是：从 $G\times G$ 限制到 $G^{\mathrm{diag}}$-等变性，再在群中作乘法。（在学习有限规范理论时，你应当已对有限群 $G$ 做过这个习题。）
\item 证明由带两个出圆周的柱面所诱导的对应图是反对角包含 $G/G\to G/G\times G/G$。
\end{enumerate}
\end{exercise}

\begin{theorem}[Freed--Hopkins--Teleman [13]，Teleman--Woodward [28]]
\begin{enumerate}[label=(\roman*),leftmargin=2.6em,itemsep=2pt]
\item 存在一个基于空间 ${}^{h+c}K^{\dim G}_G(G)$ 的二维 TQFT。
\item 带 $p$ 个输入与 $q>0$ 个输出的曲面诱导线性映射
\[
(\pi_+)_*\circ\pi_-^*\colon{}^\tau K^{\dim G}_G(G)^{\otimes p}\to{}^\tau K^{\dim G}_G(G)^{\otimes q}.
\]
\item 作为Frobenius结构基础的 ${}^\tau K^{\dim G}_G(G)^{\otimes2}$ 中的余形式由带两个出端的柱面给出，且非退化；其逆是迹配对。
\item 闭曲面的配分函数计算 $\mathcal{O}(h)$ 在 $\Phi(\Sigma;G)$ 上的指标。
\end{enumerate}
\end{theorem}

上述陈述掩盖了 [13] 中讨论的许多精细之处，例如在拉回与推前中如何处理扭转，以及 $\Phi$ 上的 K-理论定向——当 $G$ 非单连通时，后者必须小心追踪。

\subsection{推广：以 Higgs 丛为例}

该理论可以推广，以纳入 $\tilde\Phi(\Sigma;G)$ 上 Atiyah--Bott 生成上同调类的 K-理论类比。细节见 [27, 28]；线丛情形没有的一个微妙之处出现在陈述 (iv) 中：那里平坦联络的模空间必须换成所有全纯 $G_\C$-丛（包括不稳定者）的模叠。

这里我只在一个特殊情形给出 K-理论积分的公式，即在数学其他领域很有名的 Higgs 丛模空间上；积分公式的渐近推导留作读者的练习。上同调的积分公式最初由 Moore--Nekrasov--Shatashvili [21] 对 $\mathrm{SU}(n)$ 用物理的论证得到。

$\Sigma$ 上所有全纯 $G_\C$-丛的模叠 $M:=M(\Sigma;G)$ 以簇 $\Phi(\Sigma;G)$ 为相伴的 GIT 商，即
\[
\Phi(\Sigma;G)=\operatorname{Proj}\Bigl(\bigoplus_h\Gamma(M;\mathcal{O}(h))\Bigr).
\]
在一定条件下（作一个大的 $\mathcal{O}(h)$ 扭转之后），$M$ 上与 $\Phi(\Sigma;G)$ 上向量丛的指标一致。这常常使我们能在指标及相关 TQFT 的故事里免去叠，正如对线丛所做的那样。

现在 $M$ 有余切叠 $T^*M$；通常这会是微分分次叠，但在亏格 $2$ 或更高时，dg 结构消失，$T^*M$ 是普通的（局部有限 Artin）叠，局部可表为局部完全交簇被约化群作商。用同一个线丛 $\mathcal{O}(h)$，可以定义相伴的模空间 $H(\Sigma;G)$，即半稳定 Higgs 丛的模空间，它是 $T^*\Phi(\Sigma;G)$ 的部分紧化。有类似的定理把 $T^*M$ 上与 $H$ 上向量丛的指标等同起来，它特别适用于 $h>0$ 时的 $\mathcal{O}(h)$；所以不喜欢叠的人又可以在 Higgs 故事里避开叠。$H$ 的非紧性会使指标问题给出无穷的答案；但若追踪 $T^*M$ 纤维上的 $\C^*$-伸缩作用，就能使答案有限。指标将不再是数，而是 $q\in\C^*$ 的幂级数，其系数标记各权空间的维数。有了这些预备，我们就可以陈述下述定理了。

\begin{theorem}[仿照 [28]]
对 $G=\mathrm{SU}(2)$，$q$-指标 $\operatorname{Ind}(H(\Sigma;G);\mathcal{O}(h))$ 由下式给出\footnote{感谢 Daniel Halpern-Leister 指出原稿中该公式的错误。}：
\[
\begin{aligned}
&\Bigl(\frac{2h+4}{1-q}\Bigr)^{g-1}\sum_{k=1}^{h+1}\bigl(2\sin\varphi_k(q)\bigr)^{2-2g}\cdot
\bigl((1+q^2)-2q\cos(2\varphi_k(q))^{1-g}\bigr)\\
&\qquad\cdot\Bigl(1+\frac{4q}{h+2}\cdot\frac{\cos(2\varphi_k(q))-q}{(1+q^2)-2q\cos(2\varphi_k(q))}\Bigr)^{g-1},
\end{aligned}
\]
其中 $\varphi_k(q)=\pi k(q)/(h+2)$，而点 $k(q)=k+qk_1+q^2k_2+\dots$（$k=1,2,\dots,h+1$）是方程
\[
k(q)+\frac{1}{\pi i}\log\bigl(1-qe^{-2\pi i k(q)/(h+2)}\bigr)
-\frac{1}{\pi i}\log\bigl(1-qe^{2\pi i k(q)/(h+2)}\bigr)=k
\]
的幂级数解。
\end{theorem}

\begin{annote}
定理 3.2.8 主式中，底本把指数 $1-g$ 印在 $\cos(2\varphi_k(q))$ 上（即写在最后一个右括号之内）；按公式结构（与下一因子同阶、且 $q=0$ 时须退化为 Verlinde 公式），它疑应作用于整个方括号，即 $\bigl((1+q^2)-2q\cos 2\varphi_k(q)\bigr)^{1-g}$。译文照排底本。[译者按（上下文推断，无外部依据）]
\end{annote}

这是幂级数环 $\C[[q]]$ 上的一个 TQFT。在 $q=0$ 处我们得到 Verlinde 公式，理当如此，因为 $q^0$ 计数的是沿 $T^*M$ 的纤维为常值的截面，它们因而来自 $M$。对一切紧群都有推广，也可以纳入 $\mathcal{O}(h)$ 以外的向量丛；但在那种情形，我们通常必须在余切叠 $T^*M$ 上取指标，而不是在 Higgs 模空间 $H$ 上。
